\chapter*{}
\vspace{-3em}
{\raggedleft\huge\bfseries Introduction générale\par}
\vspace{2em}
\label{chap:intro}

De nos jours, la gestion de projet et la collaboration au sein des équipes constituent des enjeux essentiels pour les organisations. Face à la complexité croissante des projets, à la multiplication des tâches et à la nécessité d'assurer une bonne coordination entre les membres, l'utilisation d'outils numériques adaptés devient indispensable. Les équipes de développement logiciel, en particulier, ont besoin de solutions permettant de planifier les travaux, suivre l'avancement, gérer les priorités et favoriser une communication efficace.

Planora se présente comme une application web de gestion de projet conçue pour répondre aux besoins des équipes travaillant selon une approche collaborative et agile. Elle permet de centraliser les projets, les espaces de travail, les tâches, les backlogs, les sprints, les réunions et les échanges entre les membres. L'application intègre également des fonctionnalités intelligentes liées à l'analyse de la santé de l'équipe, afin d'aider les responsables à mieux comprendre l'état général du groupe et à détecter les signes de stress ou de frustration.

C'est dans ce cadre que s'inscrit ce projet de fin d'études, visant à concevoir et développer une plateforme complète de gestion de projet nommée Planora. Ce projet englobe la mise en place d'une architecture web basée sur un backend \ac{DOTNET}, un frontend Angular, ainsi qu'un module d'intelligence artificielle dédié à l'analyse des sentiments à partir des messages échangés entre les membres de l'équipe. L'objectif principal est de proposer une solution centralisée, ergonomique et évolutive permettant d'améliorer l'organisation, le suivi et la collaboration au sein des projets.

Le présent rapport est articulé autour des chapitres suivants :
\begin{itemize}
  \item Dans le premier chapitre, intitulé « Contexte général », nous présentons le cadre général du projet Planora, son domaine d'application, l'étude de l'existant ainsi que la méthodologie de gestion de projet retenue.
  \item Dans le deuxième chapitre, nommé « Analyse et spécification des besoins », nous identifions les acteurs du système et définissons les besoins fonctionnels et non fonctionnels. Nous y détaillons les principales fonctionnalités de l'application à travers des diagrammes de cas d'utilisation.
  \item Dans le troisième chapitre, « Conception », nous présentons l'architecture générale de la solution, la modélisation des données, les diagrammes de classes et les principaux scénarios d'interaction entre les différents composants du système.
  \item Dans le quatrième et dernier chapitre, intitulé « Réalisation », nous décrivons l'environnement de développement, les outils et technologies employés, ainsi que les interfaces les plus représentatives de l'application réalisée.
\end{itemize}

Ce rapport se clôturera par une conclusion synthétisant les acquis du projet et mettant en lumière les perspectives d'amélioration envisageables pour enrichir davantage les fonctionnalités de Planora.